% Cote 119: offprints, reports, courses, CV and letters (typed and printed part): archive cote 119, site page 66 (PDF page 67).
% offprint. Draft compared: none.
% AI-transcribed by Claude Opus 5.5 (claude-opus-5-5, adaptive thinking, effort high), Message Batches API, batch r119_full (msgbatch_01DwZNm5Y7YsKBWyAN4sFs8w), 2026-11-13; not yet checked by eye.
\scanpage{119/66}
\mspage{( 2 )}

on a $\mathrm{Z}_0 = \mathrm{Ker}\,(t : \mathrm{Z} \to \mathbf{G}_m)$ et on voit aussitôt que $\varepsilon = w\,(-1) \in \mathrm{Z}\,(\mathbf{R})$
se trouve dans le sous-groupe $\mathrm{Z}_0\,(\mathbf{R})$. On peut aussi regarder $\varepsilon$ comme
un morphisme de $\mathbf{R}$-groupes
\[ \varepsilon : \ \mu_2 \to \mathrm{Z}_0. \]

Proposition 1.3. --- La correspondance $T \to (C_0, \varepsilon)$ établit une 2-équi-%
valence de la 2-catégorie des triples de Tate avec celle des couples $(C_0, \varepsilon)$
d'une catégorie tannakienne $C_0$ (sur $\mathbf{R}$) et d'un morphisme $\varepsilon : \mu_2 \to \mathrm{Z}_0$.
1.4. On dit qu'un triple de Tate $(C, w, \mathrm{T})$ est neutre si $C$ l'est, i. e.
s'il existe un foncteur fibre $\omega : C \to \mathit{Mod\,f}\,(\mathbf{R})$. On prouve que $C$ est
neutre si et seulement si $C_0$ l'est. La donnée d'un triple de Tate neutra-%
lisé (i. e. muni d'un foncteur fibre) revient à celle d'un triple $(\mathrm{G}, w, t)$
d'un $\mathbf{R}$-groupe algébrique affine G, et des morphismes
\[ \mathbf{G}_m \xrightarrow{w} \mathrm{G} \xrightarrow{t} \mathbf{G}_m \]
vérifiant $t \circ w = -2$ et $w$ central. La catégorie tannakienne $C_0$ est la
catégorie $\mathit{Rep}_0\,(\mathrm{G}_0)$ des représentations linéaires de rang fini de $\mathrm{G}_0 = \mathrm{Ker}\,(t)$.
Le triple de Tate neutralisé se récupère à partir de $(\mathrm{G}_0, \varepsilon)$ $(\varepsilon = w\,(-1))$
par
\[ \mathrm{G} = \mathrm{Coker}\,(\tilde{\varepsilon} : \mu_2 \to \mathrm{G}_0 \times \mathbf{G}_m), \]
où
\[ \tilde{\varepsilon}\,(-1) = (\varepsilon, -1), \]
la définition de $w$, $t$ étant évidente.

2. Polarisations.

Définition 2.1. --- Soit $T = (C, w, \mathrm{T})$ un triple de Tate. Une pola-%
risation (graduée) $\pi$ de $T$ consiste en la donnée pour chaque objet V
de $C$, homogène de poids $n$, d'une classe d'équivalence $\pi\,(\mathrm{V})$ de formes
de Weil $\mathrm{V} \otimes \mathrm{V} \to \mathrm{T}^{\otimes -n}$ [(${}^1$), 1.2, 1.5] de parité $(-1)^n$ vérifiant les condi-%
tions suivantes :
PG 1. Soient V, W des objets homogènes de $C$, $\varphi \in \pi\,(\mathrm{V})$, $\psi \in \pi\,(\mathrm{W})$. Alors,
$\varphi \otimes \psi \in \pi\,(\mathrm{V} \otimes \mathrm{W})$ et si V, W sont de même poids, $\varphi \oplus \psi \in \pi\,(\mathrm{V} \oplus \mathrm{W})$.
PG 2. Le morphisme identité $\mathrm{T} \otimes \mathrm{T} \to \mathrm{T}^{\otimes 2}$ appartient à $\pi\,(\mathrm{T})$.
L'ensemble des polarisations du triple $(C, w, \mathrm{T})$ sera noté $\mathrm{Pol}\,(C, w, \mathrm{T})$.

Proposition 2.2. --- Il y a une bijection canonique
\[ \mathrm{Q} : \ \mathrm{Pol}\,(C, w, \mathrm{T}) \xrightarrow{\sim} \mathrm{Pol}_\varepsilon\,(C_0), \]
où $\varepsilon = w\,(-1)$, obtenue en associant à $\pi \in \mathrm{Pol}\,(C, w, \mathrm{T})$ l'unique $\varepsilon$-pola-%
risation sur $C_0$ telle que si $\mathrm{V} \in \mathrm{ob}\,C$ est de degré $n$, et si $\varphi \in \pi\,(\mathrm{V})$,
$\mathrm{Q}\,(\varphi) \in \mathrm{Q}\,(\pi)\,(\mathrm{Q}\,(\mathrm{V}))$.

\mspage{\hand{66}}
